\pdfvariable trailerid {[<C2762026090100000000000000000000><C2762026090100000000000000000000>]}
\documentclass[10pt]{article}
\usepackage[margin=0.68in]{geometry}
\usepackage{amsmath,amssymb,amsthm,booktabs,microtype,hyperref}
\hypersetup{colorlinks=true,linkcolor=blue,urlcolor=blue}
\linespread{0.98}
\providecommand{\CRevisionRound}{2}
\newtheorem{theorem}{Theorem}
\newtheorem{lemma}{Lemma}
\newcommand{\Pp}{\mathbb P}
\newcommand{\Ee}{\mathbb E}
\newcommand{\fall}[2]{(#1)_{#2}}
\newcommand{\ind}{\mathbf 1}
\title{Uniform Random Mapping Functional-Graph Closure:\\
Exact Cycle--Component Counts and Marked-Orbit Limits}
\author{Route-A source-local certificate HCS-C276}
\date{1 September 2026}
\begin{document}
\maketitle

\begin{abstract}
For a uniform random function $f:[n]\to[n]$, we give a self-contained
uniform random mapping functional-graph closure.  Decomposing the graph into a
permutation and a rooted forest yields the exact joint count of cyclic vertices
and components, its marginal, and every expected cycle count.
\ifnum\CRevisionRound>0
An ordered-prefix argument gives the complete tail--cycle law of a marked
orbit, identifies its collision length in distribution with the number of
cyclic vertices, and proves their common Rayleigh limit.
\fi
\ifnum\CRevisionRound>1
Conditional uniformity then gives the two-dimensional limiting density
$e^{-(x+y)^2/2}$; exhaustive independent computation audits every map through
$n=7$ without replacing the all-size proof.
\fi
\end{abstract}

\section{Model and theorem}
Choose each of the $n^n$ functions $f:[n]\to[n]$ with equal probability.
Its directed functional graph has edges $v\to f(v)$.  Let $C_n$ be the number
of cyclic vertices and $K_n$ its number of weak components.  Write
$c(k,r)$ for the unsigned Stirling number of the first kind and
$\fall nk=n(n-1)\cdots(n-k+1)$.

\begin{theorem}[Finite and asymptotic closure]\label{thm:main}
For $1\le r\le k\le n$,
\begin{equation}
 \#\{f:C_n=k,K_n=r\}
 =\binom nk c(k,r)\,k n^{n-k-1},\label{eq:joint}
\end{equation}
where the last factor means the unique empty forest when $k=n$.  Hence
\begin{equation}
 \Pp(C_n=k)=\frac{\fall nk\,k}{n^{k+1}},\qquad
 \Ee Z_{n,\ell}=\frac{\fall n\ell}{\ell n^\ell}.\label{eq:marginal}
\end{equation}
where $Z_{n,\ell}$ counts directed $\ell$-cycles.
\ifnum\CRevisionRound>0
Fix a starting vertex.  If $\mu$ is its tail length, $\lambda$ its eventual
cycle length, and $R_n=\mu+\lambda$, then
\begin{equation}
 \Pp(\mu=u,\lambda=\ell)
 =\frac{\fall{n-1}{u+\ell-1}}{n^{u+\ell}}
 \quad(u\ge0,\ \ell\ge1,\ u+\ell\le n).\label{eq:marked}
\end{equation}
$R_n\overset d=C_n$, and both variables divided by $\sqrt n$ converge to the
Rayleigh density $x e^{-x^2/2}\ind_{x\ge0}$.
\fi
\ifnum\CRevisionRound>1
Moreover,
\begin{equation}
 (\mu/\sqrt n,\lambda/\sqrt n)\Rightarrow(U,V),\qquad
 f_{U,V}(x,y)=e^{-(x+y)^2/2}\ind_{x,y\ge0}.\label{eq:jointlimit}
\end{equation}
\fi
\end{theorem}
These distributions are classical random-mapping statistics
\cite{Harris1960,FO1990}.  Our contribution is a source-locked proof synthesis
and executable boundary audit, not a literature-priority claim.

\section{Permutation and forest proof}
Fix the set $S$ of $k$ cyclic labels.  The restriction to $S$ is a
permutation, and $c(k,r)$ permutations have $r$ cycles.  Every label outside
$S$ lies in an in-tree directed toward a root in $S$.  The corresponding
complete-graph Laplacian minor, of size $m=n-k$, is
\[
 L_S=nI_m-J_m.
\]
For $m>0$ its eigenvalues are $k$ once and $n$ with multiplicity $m-1$, so
$\det L_S=k n^{n-k-1}$.  At $m=0$ the empty determinant is one.  Choosing
$S$, its permutation, and this forest is bijective with the maps in
\eqref{eq:joint}.

Summing over $r$ and using $\sum_r c(k,r)=k!$ gives
\[
 \#\{f:C_n=k\}=\binom nk k!\,k n^{n-k-1}.
\]
After division by $n^n$, this is the first formula in \eqref{eq:marginal}.
There are
$\binom n\ell(\ell-1)!=\fall n\ell/\ell$ possible directed
$\ell$-cycles.  Each prescribes $\ell$ independent function values and has
probability $n^{-\ell}$, so indicator linearity proves the second formula.
In particular, $n=1$ gives the single fixed-point map, while $k=n$ reduces
\eqref{eq:joint} to $c(n,r)$ as it must for permutations.

\ifnum\CRevisionRound>0
\section{Marked orbit and the Rayleigh bridge}
Put $t=u+\ell$.  Before its first repeat, the marked orbit visits $t$ distinct
labels.  After the fixed starting label, the ordered prefix has
$\fall{n-1}{t-1}$ choices.  Its next edge is forced to the prefix label at
index $u$, and all $n-t$ unexposed function values remain arbitrary.  The
cell therefore contains $\fall{n-1}{t-1}n^{n-t}$ maps, proving
\eqref{eq:marked}.

For each $t$ there are exactly $t$ possible tail positions.  Consequently
\begin{equation}
 \Pp(R_n=t)=\frac{t\fall{n-1}{t-1}}{n^t}
 =\frac{\fall nt\,t}{n^{t+1}}=\Pp(C_n=t).\label{eq:identity}
\end{equation}
which proves the finite distributional identity.  It also proves that
$\mu$ conditional on $R_n=t$ is uniform on $\{0,\ldots,t-1\}$.

No collision during the first $m$ transitions has probability
\begin{equation}
 \Pp(R_n>m)=\frac{\fall{n-1}m}{n^m}
 =\prod_{j=1}^m\left(1-\frac jn\right).\label{eq:survival}
\end{equation}
For $m=\lfloor x\sqrt n\rfloor$, logarithmic expansion gives
$\log\Pp(R_n>m)=-m(m+1)/(2n)+O(m^3/n^2)\to-x^2/2$.
The bound $\log(1-z)\le-z$ supplies tail tightness, so
\eqref{eq:survival} converges to the Rayleigh survival function.
Equation~\eqref{eq:identity} transfers the limit to $C_n$.
The boundary cells $\mu=0$, $\lambda=1$, and $t=n$ are already included in
\eqref{eq:marked}, without limiting conventions.
\fi

\ifnum\CRevisionRound>1
\section{The two-dimensional limit}
Let $W_n=\mu/R_n$.  Given $R_n=t$, this variable is uniform on the grid
$\{0,1/t,\ldots,(t-1)/t\}$.  The Rayleigh limit implies $R_n\to\infty$ in
probability.  Conditional Riemann sums therefore show
\[
 (R_n/\sqrt n,W_n)\Rightarrow(S,W),
\]
where $S$ has density $s e^{-s^2/2}$ and $W$ is independent uniform on
$[0,1]$.  Under $(s,w)\mapsto(sw,s(1-w))$, the Jacobian is $s$; cancellation
with the radial factor gives \eqref{eq:jointlimit}.  Its normalization is
transparent:
\[
 \int_0^\infty\!\int_0^\infty e^{-(x+y)^2/2}\,dx\,dy
 =\int_0^\infty s e^{-s^2/2}\,ds=1.
\]

\section{Executable certificate and scope}
Exact receipts enumerate all 873,612 maps for $1\le n\le7$, closing 84
cycle--component cells, 84 marked cells, and 28 cycle-length aggregates.
Formula atlases add 528 cyclic masses, 560 collision tails, and 528 cycle
expectations; 28 high-precision samples audit the scaling code.  A
producer-independent orbit tracer closes 821 assertions, SymPy closes 918
identities, fresh replay is byte-identical,
and repaired-hash hostile testing rejects 24/24 changes.  These computations
audit formulas and boundaries; Theorem~\ref{thm:main} is proved above.

Under \texttt{NO\_BAD\_EULER\_OR\_ROOT\_NUMBER}, ensemble averages supply no
arithmetic local datum, Euler factor, root number, automorphy statement, target
divisor, functional equation, or Hilbert--P\'olya operator.  There is also no
natural same-clock self-adjoint operator for this random ensemble.  The tuple
is
\[
\texttt{(A0\_FAIL,A1\_FAIL,A2\_FAIL,A3\_FAIL,A4\_FAIL)},
\]
the verdict is \texttt{ROUTE\_A\_REJECTED}, and Route B is disabled.
\fi

\begin{thebibliography}{9}
\bibitem{Harris1960}
B. Harris,
\emph{Probability distributions related to random mappings},
Ann. Math. Statist. \textbf{31} (1960), 1045--1062,
\href{https://doi.org/10.1214/aoms/1177705677}
{doi:10.1214/aoms/1177705677}.

\bibitem{FO1990}
P. Flajolet and A. M. Odlyzko,
\emph{Random mapping statistics},
Advances in Cryptology---EUROCRYPT '89, LNCS 434 (1990), 329--354,
\href{https://doi.org/10.1007/3-540-46885-4_34}
{doi:10.1007/3-540-46885-4\_34}.
\end{thebibliography}
\end{document}
